\documentclass[preprint,12pt,authoryear]{elsarticle}

\usepackage{amssymb}
\usepackage{amsmath}
\usepackage{booktabs}
\usepackage{hyperref}
\usepackage{graphicx}
\usepackage{geometry}
\geometry{a4paper, margin=1in}

\journal{Journal of Retailing and Consumer Services}

\begin{document}

\begin{frontmatter}

\title{AI-Driven Personalized Recommendation Systems and Online Impulse Buying: A Systematic Literature Review and Bibliometric Coupling Analysis of Cognitive-Affective Mechanisms}

\author[inst1]{PhD Candidate / Lab Director\corref{cor1}}
\ead{phd.director@h-ai-lab.edu.vn}
\cortext[cor1]{Corresponding author}

\author[inst1]{H-AI ResearchLAB Team}

\affiliation[inst1]{organization={Doctoral School of Business Administration, H-AI ResearchLAB},
            city={Ho Chi Minh City},
            country={Vietnam}}

\begin{abstract}
The proliferation of artificial intelligence (AI)-driven recommendation engines in contemporary electronic commerce has profoundly altered consumer browsing patterns, systematically transitioning digital shopping from intentional product retrieval into algorithmically curated discovery. This study conducts a rigorous systematic literature review following the PRISMA 2020 protocol integrated with a large-scale bibliometric coupling analysis of 1,441 peer-reviewed articles. Synthesizing these intellectual streams, we propose an integrative theoretical framework establishing cumulative e-satisfaction as the indispensable affective-cognitive bridge translating perceived recommendation accuracy and novelty into the immediate urge to buy impulsively.
\end{abstract}

\begin{keyword}
AI-driven recommendation systems \sep Online impulse buying \sep E-satisfaction \sep PRISMA 2020 \sep Bibliometric mapping \sep S-O-R framework
\end{keyword}

\end{frontmatter}

\section{Introduction}
In modern digital retailing environments, consumers are perpetually immersed in algorithmically curated information streams \citep{shanahan2019getting, xiao2018empirical}. The historical paradigm of electronic commerce has been progressively superseded by autonomous recommendation architectures powered by machine learning algorithms \citep{alnatour2011adoption, zhang2021personalized}.

\section{Methodology: The PRISMA 2020 Protocol}
Following the Preferred Reporting Items for Systematic Reviews and Meta-Analyses (PRISMA 2020) statement \citep{page2021prisma}, an initial identification query returned 2,500 bibliographic records. After systematic screening and eligibility appraisal, 1,441 clean peer-reviewed articles satisfied all methodological quality criteria, and 866 core studies with 250 links were mapped using VOSviewer.

\section{Integrative Theoretical Synthesis}
Synthesizing the evidence from the 1,441 articles, we formulate an integrated conceptual framework:
\begin{equation}
\text{Algorithmic Attributes (PP, ACC, NOV)} \longrightarrow \text{E-Satisfaction (ESAT)} \longrightarrow \text{Urge to Buy (UBI)} \longrightarrow \text{Impulse Buying (OIB)}
\end{equation}

\section{Discussion and Managerial Implications}
Retailers deploying algorithmic recommendation agents must balance recommendation accuracy and novelty to foster sustained e-satisfaction rather than triggering choice overload or surveillance anxiety.

\bibliographystyle{elsarticle-harv}
\bibliography{references}

\end{document}
